% !Mode:: "TeX:UTF-8"
% 附录。一个文件里可以写几章，每章一个 \chapter，章号换成 A、B……，图表公式跟着按
% “A-1”编号。附录 A 照范例的内容补了几样常用的；附录 B 是范例里没有、写论文常用到的
% 几种写法，照抄改成自己的即可，不用的删掉。

\chapter{补充数据}[Supplementary data]

正文里放不下的原始数据表放在这里。

\begin{table}[!ht]
  \centering
  \caption{试样的孔隙率与渗透率}[Porosity and permeability of samples]
  \label{tab:a-1}
  \begin{booktabs}{colspec={ccc}}
    \toprule
    试样 & 孔隙率 & 渗透率（$\times 10^{-12}$ m$^2$） \\
    \midrule
    1号 & 0.35 & 1.2 \\
    2号 & 0.41 & 1.8 \\
    \bottomrule
  \end{booktabs}
\end{table}

附录里的表照样按章编号，见\wref{tab:a-1}。

\section{测量不确定度}[Measurement uncertainty]

\begin{equation}\label{eq:a-1}
  u_{\mathrm{c}} = \sqrt{\sum_{i=1}^{n} c_i^2 u_i^2}
\end{equation}

公式编号也跟着换成附录的章号，见\wref{eq:a-1}。

\section{数据处理程序}[Source code]

写了题注的代码块是清单，与图表一样按章编号，见\wref{code:a-1}；加不加框、编不编行号
在 \verb|raw-style| 里改。

\begin{lstlisting}[language=Python, caption={{由压力差求渗透率}[Permeability from the pressure difference]}, label=code:a-1]
def permeability(dp, q, mu, L, A):
    """由压力差 dp、流量 q 求渗透率（达西定律）。"""
    return q * mu * L / (A * dp)
\end{lstlisting}

没写题注的是正文里的代码块，不加框、不编行号：
\begin{lstlisting}[language=Python]
k = permeability(dp, q, mu, L, A)
\end{lstlisting}

\chapter{其他写法}[Other usages]

\section{引用}[Citations]

引用单篇文献写\verb|\cite{键}|，如“空间交会控制\cite{chen1992}”。同一处引用多篇时把几个键写在一起，连续的序号自动压成区间，如\cite{lin1992,xin1999,zhao1998}。引文页码写在可选参数里，如\cite[8-13]{chen1992}。文献作为句子成分出现时写\verb|\inlinecite|，与正文排齐：由文献\inlinecite{kanamori1998}可知……

\section{定理与定义}[Theorems and definitions]

\begin{definition}[渗透率]
单位压力梯度下流体通过多孔介质的能力。
\end{definition}

\begin{theorem}[达西定律\cite{lin1992}]
低速层流时，多孔介质中的流量与压力梯度成正比。
\end{theorem}

\begin{lemma}
孔隙度越大，渗透率越大。
\end{lemma}

\begin{remark}
定理类环境的名称、编号与字体由模板给出，引理、推论、假设、猜想等同样可用，见手册。
\end{remark}

\section{并排的图}[Figures side by side]

两张图并排、各有各的图题时，各装一个 minipage。两图取同样的高，minipage 按首行对齐，图题才排在同一高度：

\begin{figure}[!ht]
  \centering
  \begin{minipage}[t]{0.45\textwidth}
    \centering
    \includegraphics[height=2.5cm]{aerostatic}
    \caption{气体静压轴承}[Externally pressurized gas bearing]
  \end{minipage}%
  \begin{minipage}[t]{0.45\textwidth}
    \centering
    \includegraphics[height=2.5cm]{aerodynamic}
    \caption{气体动压轴承}[Aerodynamic gas bearing]
  \end{minipage}
\end{figure}

分图题要中英文两行时，把\verb|\bisubcaption{中文}{English}|写进\verb|\subfigure|的可选参数。

\section{算法}[Algorithms]

规范里没有算法的格式。文档类装好了 \verb|algorithm2e|，算法按章编号，题注与图表一样写，英文题写在方括号里；外框、关键字大小写、缩进线、行号的样子在 \verb|style/algorithm| 里挑，见手册：

\begin{algorithm}[!ht]
  \caption{渗透率的迭代求解}[Iterative solution of permeability]
  \KwIn{压力差 $\Delta P$，流量 $M$}
  \KwOut{渗透率 $\alpha$}
  $\alpha \gets \alpha_0$\;
  \While{$|\alpha - \alpha'| > \varepsilon$}{
    $\alpha' \gets \alpha$\;
    由\wref{eq:2-1}更新 $\alpha$\;
  }
  \Return{$\alpha$}
\end{algorithm}

\section{脚注}[Footnotes]

脚注\footnote{脚注的字号、行距与编号样式由模板给出。}写\verb|\footnote|，编号是带圈数字，每页重新编号。
